\documentclass{article}
\usepackage[T2A]{fontenc}
\usepackage[cp1251]{inputenc}
\usepackage{amsthm}
\usepackage{amsmath}
\usepackage{amssymb}
\usepackage{amsfonts}
\usepackage{mathrsfs}
\usepackage[12pt]{extsizes}
\usepackage{fancyvrb}
\usepackage{indentfirst}
\usepackage[
  left=2cm, right=2cm, top=2cm, bottom=2cm, headsep=0.2cm, footskip=0.6cm, bindingoffset=0cm
]{geometry}
\usepackage[english,russian]{babel}
\begin{document}
\section*{Вариант 9}
\begin{equation}
R^{*}_{A_j}(\omega) =
\begin{cases}
(1 - (t_0\omega)^2) / (1 + (t_0\omega)^4), & A_j(\mathbf{p}, \mathbf{s}) \neq 0, \\
\sqrt{1 + \omega^2}(1 - (t_0\omega)^2) / (1 + (t_0\omega)^4), & A_j(\mathbf{p}, \mathbf{s}) = 0,
\end{cases}
\end{equation}
\begin{equation*}
A_j(\mathbf{p}, \mathbf{s}) = \left[\sum\limits_{\nu=1}^{N_y} |\Phi_{\nu j}(0, \mathbf{p}, \mathbf{s})|^2\right]^{1/2},
\end{equation*}
\begin{equation*}
\Psi_{\nu j}(\lambda, \mathbf{p}, \mathbf{s}) = \Phi_{\nu j}(\lambda, \mathbf{p}, \mathbf{s}) / \lambda,
\quad \nu = 1, 2, \dots, N_y; \quad j = 1, 2, \dots, N_x,
\end{equation*}
где $\mathbf{p} \in \Omega^{(st)}_p(\mathbf{s}_0)$~--- набор параметров обратных связей в момент
старта параметрического синтеза, $t_0$~--- желаемое время регулирования.
Поскольку функция $F : \mathbb{R}^{N_p} \to \mathbb{R}$ является негладкой,
а размерность $N_p$ пространства параметров обратных связей обычно не превышает
нескольких десятков, для минимизации обычно используется безградиентный метод
Нелдера"---Мида.
\end{document}